\section*{Abstract}

Dyslexia screening in low-resource languages is constrained by the scarcity of expert clinicians, validated test batteries and labelled corpora. Multilingual large language models (LLMs) are an attractive substrate because their pretraining already encodes substantial cross-lingual structure, but it is not known which architectures actually carry lexical and phonological information that transfers from a high-resource source (English) to a low-resource target (Bangla), nor where in the network such information lives. This thesis probes four publicly released multilingual LLMs that span the four dominant pretraining recipes -- XLM-RoBERTa large (multilingual masked encoder, 0.56B), Llama 3.1 8B (English-heavy causal decoder), BLOOM 7B1 (46-language causal decoder) and mT5 XL (multilingual span-corruption encoder-decoder, 3.7B). A curated set of 440 stimuli (220 English + 220 Bangla) spanning seven phonological task families and a cross-lingual concept-matching set is fed through each frozen model, and per-layer linear probes are trained for voicing, real-versus-pseudo word, orthographic transparency and syllable count.

A central methodological finding is that three of these four standard tasks fail surface-feature sanity checks: voicing labels are recoverable from the first character alone, orthographic and syllable probes do not improve over a five-feature character-statistic baseline, and the apparent cross-lingual signal in prior work largely reflects character statistics or token identity rather than phonological structure. After length-matched stimulus reconstruction, only real-versus-pseudo word discrimination survives. On this surviving task, BLOOM reaches 0.92 F1 on Bangla zero-shot transfer and XLM-R reaches 0.86 F1 with only 0.56B parameters, while Llama (0.63 F1) and mT5 (0.60 F1) collapse to surface-baseline level for cross-lingual transfer. Representational Similarity Analysis (RSA) and causal forward-hook ablation confirm that BLOOM and XLM-R distribute cross-lingual lexical information widely, whereas mT5 alone localises it: zeroing ten ranked units in the encoder selectively impairs Bangla transfer by 0.11--0.14 F1 across all probed encoder depths, and two units (782 and 1880) persist as a backbone across more than half of those depths. The thesis concludes with a paired recommendation -- BLOOM 7B1 for practical low-resource Bangla screening prototypes and mT5 XL for interpretable, mechanism-revealing analyses -- and with a methodology contribution: surface-matched stimuli plus native-verified labels are the minimum standard for any future clinical probing of multilingual LLMs.

\vspace{1cm}
\textbf{Keywords:} dyslexia screening; multilingual large language models; cross-lingual transfer; probing; representational similarity analysis; causal ablation; low-resource Bangla NLP; mechanistic interpretability.
\pagebreak
